\section{Conclusion}

\subsection{Résumé des travaux réalisés}

Ce projet de recherche s'est attaché à l'élaboration d'une solution algorithmique de segmentation sémantique appliquée à l'imagerie ultrasonore, dans le contexte spécifique de l'analyse de puits de forage. L'enjeu principal résidait dans l'identification précise de trois structures géologiques et mécaniques distinctes (Ciment, Casing, TIE) au sein d'un jeu de données restreint et présentant un fort déséquilibre de classes.

Notre méthodologie a reposé sur une exploration systématique de différentes stratégies de modélisation :

\begin{itemize}
      \item \textbf{Architecture de référence} : Implémentation d'un réseau U-Net standard à initialisation aléatoire (31M paramètres).
      \item \textbf{Variantes architecturales} : Intégration de mécanismes de régularisation (Dropout) et d'attention (Attention Gates).
      \item \textbf{Enrichissement des données} : Application de techniques d'augmentation (transformations géométriques, injection de bruit gaussien).
      \item \textbf{Fonctions de coût} : Évaluation comparative de critères d'optimisation (Focal Loss, Dice Loss, Combined Loss).
      \item \textbf{Apprentissage par transfert} : Exploitation d'encodeurs pré-entraînés (famille ResNet).
      \item \textbf{Stratégies d'entraînement} : Mise en œuvre de protocoles de \textit{fine-tuning} progressif.
\end{itemize}

\subsection{Conclusion et résultats}

Ce travail de recherche a permis d'optimiser une architecture de segmentation pour l'imagerie de puits en environnement contraint (données limitées, déséquilibre de classes). L'analyse comparative consacre le modèle U-Net avec encodeur ResNet-34 pré-entraîné et fonction de coût \textit{Dice Loss} comme solution optimale, validant la supériorité du transfert d'apprentissage pour cette tâche.

Avec un protocole d'entraînement séquentiel et des augmentations par bruit gaussien, cette configuration atteint une \textbf{IoU de 0.885} sur l'ensemble de validation (+8.95\% par rapport à la baseline). Ce résultat confirme qu'un modèle à forte capacité de représentation, correctement régularisé, offre le meilleur compromis performance-généralisation.

Finalement, ces performances se traduisent par un classement de 15\textsuperscript{e}/52 au niveau global (4\textsuperscript{e} de la promotion), soulignant la validité d'une démarche méthodique combinant connaissances métiers et stratégies avancées de Deep Learning.